\chapter{Restfolgen, Teilfolgen und Indexverhalten}
\label{b21guide:ch:subsequences}

Hier ist $a\colon\mathbb N\to M$ weiterhin eine Folge in einer
beliebigen Menge $M$. Konvergenz wird noch nicht vorausgesetzt.

\section{Schließlich und unendlich oft}

\begin{definition}[Aussagen über späte Indizes]
\label{b21guide:def:eventually}
Für eine Aussage $P(n)$ über $n\in\mathbb N$ bedeutet
\[
 \begin{aligned}
 P(n)\text{ gilt schließlich}
 &\quad\Longleftrightarrow\quad
 \exists N\in\mathbb N\;\forall n\geq N:\ P(n),\\
 P(n)\text{ gilt unendlich oft}
 &\quad\Longleftrightarrow\quad
 \forall N\in\mathbb N\;\exists n\geq N:\ P(n).
 \end{aligned}
\]
\end{definition}

\begin{theorem}[Endliche Ausnahmen und Negation]
Eine Aussage gilt schließlich genau dann, wenn sie an höchstens
endlich vielen Indizes falsch ist. Sie gilt unendlich oft genau
dann, wenn die Menge ihrer gültigen Indizes unendlich ist. Ferner
\[
 \begin{aligned}
 \neg(P(n)\text{ schließlich})
 &\ \Longleftrightarrow\ \neg P(n)\text{ unendlich oft},\\
 \neg(P(n)\text{ unendlich oft})
 &\ \Longleftrightarrow\ \neg P(n)\text{ schließlich}.
 \end{aligned}
\]
Endlich viele schließlich gültige Aussagen gelten schließlich
gleichzeitig.
\end{theorem}

\section{Restfolgen und Teilfolgen}

\begin{definition}[Restfolge und Teilfolge]
\label{b21guide:def:subsequence}
Für $r\in\mathbb N$ heißt
$a^{\langle r\rangle}_n=a_{n+r}$ die \emph{Restfolge ab $r$}.
Eine \emph{Teilfolge} von $a$ ist eine Folge
$b_n=a_{\varphi(n)}$, wobei
$\varphi\colon\mathbb N\to\mathbb N$ strikt wachsend ist:
$\varphi(n+1)>\varphi(n)$ für alle $n$.
\end{definition}

\begin{theorem}[Indizes einer Teilfolge]
\label{b21guide:thm:indices}
Für jede Teilfolgenabbildung $\varphi$ gelten
\[
 \varphi(n)\geq n,\qquad
 \varphi(n+r)\geq\varphi(n)+r
 \quad(n,r\in\mathbb N).
\]
Insbesondere überschreiten ihre Werte schließlich jede vorgegebene
natürliche Zahl. Eine schließlich für $a_n$ gültige Eigenschaft
gilt auch schließlich für $a_{\varphi(n)}$.
\end{theorem}

\begin{theorem}[Auswahl unendlich vieler Indizes]
\label{b21guide:thm:enumeration}
Jede unendliche Teilmenge $J\subseteq\mathbb N$ besitzt genau
eine strikt wachsende Aufzählung
$j_0<j_1<j_2<\cdots$, in der jedes Element von $J$ vorkommt.
Gilt eine Eigenschaft von $a_n$ unendlich oft, so gibt es daher
eine Teilfolge, deren sämtliche Glieder diese Eigenschaft haben.
\end{theorem}

\begin{theorem}[Wiederholtes Auswählen und Verschieben]
Jede Restfolge ist eine Teilfolge. Eine Teilfolge einer Teilfolge
ist wieder eine Teilfolge: Sind $\varphi,\psi$ strikt wachsend,
so hat sie die Form
\[
 \bigl(a_{\varphi(\psi(n))}\bigr)_{n\in\mathbb N}.
\]
Für $r,s\in\mathbb N$ gilt außerdem
\[
 \bigl(a^{\langle r\rangle}\bigr)^{\langle s\rangle}
 =a^{\langle r+s\rangle}.
\]
\end{theorem}

Die strikte Zunahme betrifft die \emph{Indizes}, nicht die
Folgenglieder. Eine Teilfolge darf beliebig viele gleiche Werte
enthalten.

\section{Periodische und rekursiv erzeugte Folgen}

\begin{definition}[Periodizität]
\label{b21guide:def:periodic}
Eine Folge $a\colon\mathbb N\to M$ heißt \emph{periodisch}, wenn es
eine ganze Zahl $p\geq1$ mit $a_{n+p}=a_n$ für alle $n\in\mathbb N$
gibt. Die Zahl $p$ heißt eine \emph{Periode}.
Gilt diese Gleichung nur für alle $n\geq r$ mit einem festen
$r\in\mathbb N$, so heißt die Folge \emph{schließlich periodisch}.
Eine kleinste Periode wird dabei nicht vorausgesetzt.
\end{definition}


\begin{theorem}[Kleinste Periode und gemeinsame Perioden]
\label{b21guide:thm:least-period}
Jede periodische Folge $a\colon\mathbb N\to M$ besitzt eine
kleinste positive Periode $p_0$. Für jede ganze Zahl $p\geq1$ gilt:
\[
 p\text{ ist Periode von }a
 \quad\Longleftrightarrow\quad
 p=kp_0\text{ für eine ganze Zahl }k\geq1.
\]
Sind $p,q$ Perioden, so ist auch $\operatorname{ggT}(p,q)$
eine Periode; $\operatorname{ggT}$ bezeichnet den größten
gemeinsamen Teiler der beiden positiven ganzen Zahlen.

Entsprechend besitzt jede schließlich periodische Folge eine
kleinste schließlich gültige positive Periode. Ihre schließlich
gültigen Perioden sind genau deren positive Vielfache;
auch ihr größter gemeinsamer Teiler ist wieder eine schließlich
gültige Periode. Der Beginn der Periodizität darf von der
betrachteten Periode abhängen.
\end{theorem}

\begin{theorem}[Wiederholungen bei einer festen Rekursionsvorschrift]
\label{b21guide:thm:recursive-period}
Sei $T\colon M\to M$ und $a_{n+1}=T(a_n)$ für alle $n$.
Gilt $a_r=a_s$ für zwei Indizes $r<s$, so gilt
\[
 a_{r+k}=a_{s+k}\qquad(k\in\mathbb N).
\]
Die Folge ist dann ab $r$ periodisch mit Periode $s-r$.
Hat eine solche rekursiv erzeugte Folge nur endlich viele
verschiedene Werte, so ist sie schließlich periodisch.
Die Schrittvorschrift $T$ muß hierbei unabhängig vom Index sein.
\end{theorem}
